\section{Improvements III: Modelling}
\label{sec:impmodel}

\subsection{Kernel bandwidth $\alpha$ (I3)}
\paragraph{Problem.} With $n$ qubits and inputs spread over $[0,\pi]$, the ZZ feature states of two
different flows are nearly orthogonal. Roughly, each qubit and each pair contributes an
independent ``random'' phase, and the overlap of two random states in dimension $2^n$ is of order
$2^{-n}$. The experiments confirm this: at 16 qubits with $\alpha=1$ and 2 repetitions, the
\emph{median} off-diagonal training kernel entry is about $2\times10^{-5}$ and 84--89\% of entries
are below $10^{-3}$. The mean is higher (about 0.03, Table~\ref{tab:bandwidth}) only because a
few near-identical flows still overlap strongly. An SVM can still
\emph{separate} the training set with such a kernel, but each test prediction is then decided by
a handful of nearest training points, and the kernel's geometry carries little information.

\paragraph{Change.} All feature maps now take a bandwidth parameter. Angles from preprocessing
are multiplied by $\alpha$ before encoding:
\begin{equation}
\ket{\phi_\alpha(\bm\theta)}=U(\alpha\bm\theta)\ket{0^n},\qquad \bm\theta\in[0,\pi]^n .
\end{equation}
Small $\alpha$ moves all states towards $U(\bm 0)\ket{0^n}$, raising off-diagonal kernel values.
In the limit $\alpha\to0$ the kernel tends to the constant 1. For the separable $R_x$ map,
$\alpha$ is exactly an RBF-like length-scale: $\cos^2(\alpha\Delta/2)\approx e^{-\alpha^2\Delta^2/4}$
for small $\alpha\Delta$. For the ZZ map, the pair term
$(\pi-\alpha\theta_i)(\pi-\alpha\theta_j)=\pi^2-\alpha\pi(\theta_i+\theta_j)+\alpha^2\theta_i\theta_j$
shows that $\alpha$ also controls how strongly the \emph{product} (entangling) terms depend on the
data, relative to the linear ones. $\alpha=1$ reproduces the original circuit exactly
(\code{ZZFeatureMap} and the fast engine both take \code{alpha}).

\subsection{Cross-validated tuning with equal budgets (I2)}
\paragraph{Problem.} With $C=1$ and a fixed scale for both models, the comparison measures how
well two \emph{default settings} happen to suit this dataset. RBF-SVM accuracy is very sensitive to
$\gamma$ and $C$. The default \code{gamma='scale'} ($1/(d\cdot\mathrm{Var}(X))$) is a heuristic, not
an optimum.

\paragraph{Change.} \code{evaluation.tune\_quantum()} performs a grid search over
$\alpha\in\{0.05,0.1,0.2,0.4,1.0\}$ and $C\in\{0.1,1,10,100\}$, scoring \emph{macro-F1} by 5-fold
stratified CV \emph{on the training split only}. The expensive part, the kernel, is computed once
per $\alpha$. Every $C$ and every fold then reuses sub-blocks $K[\text{tr},\text{tr}]$ and
$K[\text{va},\text{tr}]$ of that one matrix, so the search costs 5 kernel computations rather than
100 fits' worth. \code{tune\_rbf()} runs \code{GridSearchCV} for the RBF-SVM over the same
$C$ grid and $\gamma\in\{\text{scale},0.01,0.03,0.1,0.3,1\}$, with the same folds, metric and seed.
The best configuration is refitted on the full training split and evaluated once on the test set.
Both models get 20--24 candidate configurations: an equal budget.

\begin{verbatim}
alpha, C, K_train, table = tune_quantum(
    lambda a: make_kernel(a).compute_kernel_matrix(A_train, verbose=False),
    y_train, args.alphas, args.Cs, seed=seed)
\end{verbatim}

\code{--no-tune --alpha A --C C} fixes the values instead. \code{--alphas} and \code{--Cs} change
the grids.

\subsection{Entanglement topology and depth}
\code{entangling\_pairs(n, entanglement)} generalises the hard-coded adjacent pairs:
\code{linear} $\{(i,i{+}1)\}$ (original, $n-1$ pairs), \code{circular} (adds $(n{-}1,0)$), and
\code{full} (all $\binom n2$ pairs, the default in Qiskit's \code{ZZFeatureMap}). Both the Cirq
circuit builder and the fast engine use the same pair list, and \code{--n-reps} sets the depth. These
options answer the paper's ``higher-order entanglement'' future-work item within the existing
circuit family.

\subsection{All three paper methods in one pipeline (I10)}
\code{--kernel rx} selects the separable $R_x$ map (Method 1), and \code{--sampling --shots s} adds
finite-shot estimation to any kernel (Method 2). Together with the default ZZ map (Method 3), all
methods now share one loader, one split per seed, one preprocessing object and one tuning
protocol. The paper's comparisons used different $N$ for different methods (822 vs.\ 1028 binary
test samples); here every method sees exactly the same flows.

\subsection{Stronger classical references}
Besides the tuned RBF-SVM on the \emph{same} PCA features, which isolates the effect of the kernel,
\code{--extra-baselines} trains a Random Forest~\cite{breiman2001} (300 trees) on \emph{all}
raw features without PCA. This is a typical practical NIDS baseline. It answers a different
question: how much accuracy is left on the table by compressing to $n$ components at all?
